\chapter*{Conclusion générale}
\phantomsection
\addcontentsline{toc}{chapter}{Conclusion générale}

Ce projet a étudié la gravité des blessures parmi les usagers des accidents corporels
enregistrés en France métropolitaine entre 2019 et 2024. Il a suivi la chaîne complète
d'une étude de données~: compréhension des sources, audit, préparation, analyse
statistique, modélisation et évaluation sur une année ultérieure.

Le premier résultat est une table d'analyse traçable~: 702\,823 observations d'usagers
pour 309\,756 accidents distincts. Sa construction a demandé de résoudre les différences
de schéma, les codes spéciaux, les identifiants, les valeurs atypiques et les lieux
multiples. Les règles retenues préservent les usagers, signalent les désaccords et
maintiennent les valeurs sources nécessaires à la vérification. Les 365 gravités
inconnues restent identifiées, puis sont exclues des calculs supervisés.

Les analyses de 2019-2022 répondent à la question des différences observées. Les
proportions de décès ou d'hospitalisations sont plus élevées chez les piétons que chez
les conducteurs, chez les personnes de 75 ans ou plus que chez les 35-44 ans, de nuit
sans éclairage public qu'en plein jour, et hors agglomération qu'en agglomération.
Le bootstrap par accident rend compte d'une variabilité en respectant le regroupement
des usagers. La lecture annuelle et la sensibilité aux décès seuls apportent des
compléments, sans effacer les limites de collecte ni transformer ces associations
en effets causaux.

La partie prédictive apporte une réponse complémentaire. Avec onze variables et une
régression logistique régularisée, les probabilités sont mieux estimées que par la
seule fréquence de l'apprentissage. Le réglage choisi sur 2023 se maintient sur
2024~: la log-loss diminue d'environ 14,93\,\% par rapport à la référence, avec une
ROC-AUC proche de 0,770. Le projet établit ainsi un apport mesurable dans le cadre
de l'évaluation retenue.

Les diagnostics précisent la portée de ce résultat. Au seuil de 0,5, le rappel est
limité à 18,22\,\%. Les probabilités élevées sont souvent surestimées et le classement
est moins performant chez les piétons. L'absence d'âge ou de sexe porte également un
signal de saisie particulier. Le modèle conserve un gain sur les usagers dont ces
informations sont présentes, mais la stabilité de ce signal en cas de changement
de collecte n'est pas établie.

La principale contribution du projet tient à la cohérence entre la question,
l'unité d'analyse, les transformations et l'évaluation. Les résultats restent
conditionnels à la présence dans BAAC. Ils ne décrivent pas un risque par trajet,
ne permettent pas une interprétation causale des coefficients et ne démontrent
pas l'intérêt d'un outil de décision en temps réel. Cette délimitation permet
de défendre précisément ce qui a été réalisé.

Plusieurs prolongements seraient possibles~: approfondir la calibration, étudier les
mécanismes d'information manquante, représenter des interactions ou comparer une
autre famille de modèles. Ils demanderaient un protocole d'évaluation supplémentaire,
car les scores de 2024 sont désormais connus. Pour mesurer un risque par déplacement
ou éclairer une décision d'aménagement, des données d'exposition et un objectif
d'usage seraient également nécessaires.

